% Chapter 5 of Singular Learning Theory for the Rest of Us: blow-ups and the resolution of singularities.
\sltchapter{Blow-Ups}{Resolution of singularities by hand}{$y=ax_1+b^2x_2$}{ch:blowups}
\newtonfoot{0/4,2/0}{4/3}{\tfrac34}{a^2+b^4}
\begin{refsection}[chapters/05-blowups.bib]

\begin{chaptersummary}
The quartic model's loss $\frac12(a^2+b^4)$ has a Hessian of rank one at the truth: counting parameters predicts $\lambda=1$, counting nonzero Hessian eigenvalues $\frac12$, and the truth is $\frac34$. We compute it by hand with two blow-ups, changes of coordinates that replace the singular point by a curve of directions until the loss is a monomial; the cusp $(a^2-b^3)^2$ needs three and gives $\frac5{12}$. We prove that every resolution computes the volume and the zeta function, so that $\lambda$ and $m$ are read off its exponents, and the rules of the trade, among them $\lambda=\frac{d-s}2$ near a smooth $s$-dimensional set of true parameters across which the loss grows quadratically. Measured volumes agree with the charts to six digits and approach the predicted constants, slowly for the cusp. Hironaka's theorem guarantees a resolution for every analytic loss, but finding one is another matter: for most networks none is known.
\end{chaptersummary}

\section{The question: what does a flat direction cost?}\label{sec:blowups:question}

The Hessian of a trained network typically has a bulk of eigenvalues near zero and a few large ones \parencite{SagunEtAl2016},\sidenote{\textcite{SagunEtAl2016} is a preprint; \cref{ch:llc} returns to Hessian-based measures of complexity.} and a tempting reading is that the effective number of parameters is the rank of the Hessian. The model of this chapter tests that reading with two parameters. Fit $y=ax_1+b^2x_2+\varepsilon$, with $x\sim\Normal(0,I_2)$ and $\varepsilon\sim\Normal(0,1)$, to pure noise, with the uniform prior on $[-1,1]^2$. The loss of this \term{quartic model}\index{quartic model} relative to the truth $a=b=0$ is
\[
K(a,b)=\tfrac12\E_x\bigl[(ax_1+b^2x_2)^2\bigr]=\tfrac12\bigl(a^2+b^4\bigr),
\]
whose Hessian at the truth is $\diag(1,0)$: along $b$ the loss is flat to second order. Counting parameters gives $\lambda=\frac d2=1$, and counting nonzero Hessian eigenvalues gives $\frac12$. Neither is right: the free energy of the quartic model grows like $\frac34\log n$, with measured slope $0.749$ (\cref{ch:freeenergy}). The flat direction is not free, because the loss grows like $b^4$ along it; it costs a quarter where a regular direction costs a half.

\Cref{ch:volume} explained where such numbers come from when $K$ is a \term{monomial}: if $K=\prod_jw_j^{2k_j}$ near $0$ and the prior behaves like $\prod_j\abs{w_j}^{h_j}$, each coordinate with $k_j>0$ offers the candidate $(h_j+1)/(2k_j)$, the learning coefficient is the smallest candidate, and its multiplicity is the number of candidates that tie (\cref{prop:volume:monomial}). The product model's $a^2b^2$ has candidates $\frac12$ and $\frac12$, so $\lambda=\frac12$ and $m=2$. The quartic $a^2+b^4$ is not a monomial, and neither is the loss of a network. The regular theory handles one kind of non-monomial loss.

\begin{regularanswer}[the Morse lemma]
If $K\ge0$ is analytic, $K(w_0)=0$ at an interior point $w_0$, and $\nabla^2K(w_0)$ is positive definite, then there are analytic coordinates $u$ centred at $w_0$ in which $K=u_1^2+\dots+u_d^2$ (\cref{thm:toolkit:splitting} with $r=d$). Then $\{K<\eps\}$ is a ball of radius $\sqrt\eps$, $V(\eps)\sim c\,\eps^{d/2}$, and $\lambda=\frac d2$, $m=1$ for a prior positive at $w_0$ \parencite[Lemma~3]{Watanabe2013}. Every parameter costs one half.
\end{regularanswer}

The quartic model's Hessian is degenerate, and cleverer coordinates cannot help. A function is \term{normal crossing}\index{normal crossing} at a point if, in some analytic coordinates centred there, it is a monomial times an analytic function that does not vanish at the point; near the point it then vanishes exactly on a union of coordinate hyperplanes. The function $a^2+b^4$ vanishes only at the origin, and a change of coordinates that is invertible near the origin cannot turn a point into a union of lines (nor can it for the regular $\sum_iu_i^2$, $d\ge2$). The way out is a change of variables that replaces the singular point by a whole curve and is invertible everywhere else. Integrals do not notice it, because the curve is a null set; the loss does, because it can become a monomial. Such changes are \term{blow-ups}\index{blow-up}. The question of this chapter is how to find them for a given loss, and why the exponents they produce are the learning coefficient and its multiplicity.

\section{By hand: blowing up a point}\label{sec:blowups:byhand}

\paragraph{The blow-up of the plane.} Every point $w\ne0$ of the plane lies on exactly one line through the origin, and the origin lies on all of them. The \term{blow-up} of the origin separates the lines there: its points are the pairs $(w,\ell)$ of a line $\ell$ through the origin and a point $w\in\ell$, and the map $g(w,\ell)=w$ forgets the line. Away from $w=0$ the line is determined by $w$, so $g$ is one-to-one there; over the origin sits a whole curve of points $(0,\ell)$, one for each direction, the \term{exceptional divisor}\index{exceptional divisor} $E$.\sidenote{``Divisor'' is the geometers' word for a hypersurface, here a curve. The directions through the origin form the real projective line, a circle, and near $E$ the blow-up is a Möbius band.} Two charts describe it, chart I by the slope $v$ of the line $b=va$ and the abscissa $u$, chart II by $a=ub$ and the ordinate $v$:
\begin{equation}\label{eq:blowups:charts}
\text{I:}\ (a,b)=(u,uv),\ \det g'=u;\qquad\text{II:}\ (a,b)=(uv,v),\ \det g'=v.
\end{equation}
The exceptional divisor is the axis $\{u=0\}$ in chart I, whose point $(0,v)$ is the direction of slope $v$, and the axis $\{v=0\}$ in chart II. A line $b=ca$ becomes the line $\{v=c\}$ of chart I, so lines that met at the origin no longer meet (\cref{fig:blowups:plane}); where both charts apply, $(u,v)$ in chart I is $(1/v,uv)$ in chart II. The map $g$ is analytic, proper (preimages of compact sets are compact), and an analytic isomorphism off the null sets $E$ and $\{0\}$, so integrals can be computed on the charts. The part $\{\abs b\le\abs a\le1\}$ of the square is the image of $[-1,1]^2$ in chart I, and the part $\{\abs a\le\abs b\le1\}$ that in chart II, so for every integrable $F$
\begin{equation}\label{eq:blowups:twocharts}
\int_{[-1,1]^2}F(a,b)\dd a\dd b=\int_{[-1,1]^2}\bigl(F(u,uv)\abs u+F(uv,v)\abs v\bigr)\dd u\dd v.
\end{equation}
In $\R^d$ the blow-up of the origin has $d$ charts, $w_j=u_j$ and $w_i=u_ju_i$ for $i\ne j$, with Jacobian $\abs{u_j}^{d-1}$.

\begin{figure}[t]
\begin{fullwidth}
\centering
\begin{tikzpicture}
\begin{groupplot}[group style={group size=3 by 1,horizontal sep=10mm},
  width=0.33\linewidth,height=0.31\linewidth,axis lines=middle,axis line style={muted!70,-},
  xtick=\empty,ytick=\empty,clip=true,title style={font=\small,yshift=0.6ex},
  xlabel style={anchor=north west,font=\small,inner sep=1pt},ylabel style={anchor=south east,font=\small,inner sep=1pt}]
\nextgroupplot[title={chart I: $(a,b)=(u,uv)$},xlabel={$u$},ylabel={$v$},xmin=-1.15,xmax=1.15,ymin=-2.3,ymax=2.3]
\addplot[name path=IR,draw=singular!55,line width=0.5pt,domain=-2.3:2.3,samples=61,variable=\t]
  ({sqrt(0.1/(1+sqrt(1+0.2*\t^4)))},{\t});
\addplot[name path=IL,draw=singular!55,line width=0.5pt,domain=-2.3:2.3,samples=61,variable=\t]
  ({-sqrt(0.1/(1+sqrt(1+0.2*\t^4)))},{\t});
\addplot[singular!14] fill between[of=IR and IL];
\pgfplotsinvokeforeach{-2,-1,-0.5,0,1,2}{\addplot[regular!45,line width=0.5pt,domain=-1.15:1.15,samples=2] {#1};}
\addplot[regular,line width=1.1pt,domain=-1.15:1.15,samples=2] {0.5};
\node[font=\scriptsize,text=regular,anchor=south] at (axis cs:0.62,0.5) {$v=\frac12$};
\addplot[singular,line width=1.6pt] coordinates {(0,-2.3) (0,2.3)};
\node[font=\small,text=singular,anchor=west] at (axis cs:0.05,1.75) {$E$};
\nextgroupplot[title={the plane $(a,b)$},xlabel={$a$},ylabel={$b$},xmin=-1.1,xmax=1.1,ymin=-1.1,ymax=1.1]
\addplot[name path=OR,draw=singular!55,line width=0.5pt,domain=-0.4728:0.4728,samples=61,variable=\t]
  ({sqrt(max(0,0.05-\t^4))},{\t});
\addplot[name path=OL,draw=singular!55,line width=0.5pt,domain=-0.4728:0.4728,samples=61,variable=\t]
  ({-sqrt(max(0,0.05-\t^4))},{\t});
\addplot[singular!14] fill between[of=OR and OL];
\pgfplotsinvokeforeach{-2,-1,-0.5,0,1,2}{\addplot[regular!45,line width=0.5pt,domain=-1.1:1.1,samples=2] {#1*x};}
\addplot[regular!45,line width=0.5pt] coordinates {(0,-1.1) (0,1.1)};
\addplot[regular,line width=1.1pt,domain=-1.1:1.1,samples=2] {0.5*x};
\node[font=\scriptsize,text=regular,anchor=north west] at (axis cs:0.55,0.27) {$b=\frac a2$};
\fill[singular] (axis cs:0,0) circle[radius=1.9pt];
\nextgroupplot[title={chart II: $(a,b)=(uv,v)$},xlabel={$u$},ylabel={$v$},xmin=-2.3,xmax=2.3,ymin=-1.15,ymax=1.15]
\addplot[name path=JU,draw=singular!55,line width=0.5pt,domain=-2.3:2.3,samples=81,variable=\t]
  ({\t},{sqrt(0.1/(\t^2+sqrt(\t^4+0.2)))});
\addplot[name path=JD,draw=singular!55,line width=0.5pt,domain=-2.3:2.3,samples=81,variable=\t]
  ({\t},{-sqrt(0.1/(\t^2+sqrt(\t^4+0.2)))});
\addplot[singular!14] fill between[of=JU and JD];
\pgfplotsinvokeforeach{-2,-1,-0.5,0,0.5,1}{\addplot[regular!45,line width=0.5pt] coordinates {(#1,-1.15) (#1,1.15)};}
\addplot[regular,line width=1.1pt] coordinates {(2,-1.15) (2,1.15)};
\node[font=\scriptsize,text=regular,anchor=south east] at (axis cs:1.95,0.62) {$u=2$};
\addplot[singular,line width=1.6pt] coordinates {(-2.3,0) (2.3,0)};
\node[font=\small,text=singular,anchor=south] at (axis cs:-1.75,0.04) {$E$};
\end{groupplot}
\end{tikzpicture}
\end{fullwidth}
\caption{The blow-up of the origin of the plane. Middle: eight lines through the origin and the set $\{a^2+b^4<0.05\}$ (shaded). Left and right: the two charts \eqref{eq:blowups:charts}, which show all of the lines but the vertical one (chart I) and the horizontal one (chart II). Each line through the origin becomes a line crossing the exceptional divisor $E$ (bronze) once, at its direction: $b=a/2$ (dark slate) is $\{v=\frac12\}$ in chart I and $\{u=2\}$ in chart II, and the origin has become all of $E$. In chart I the shaded set is a strip of nearly constant width around $E$, because $a^2+b^4=u^2\times$\,(a unit) there; in chart II it still pinches at the centre, where $a^2+b^4=v^2(u^2+v^2)$ is not normal crossing (\cref{comp:blowups:quartic}).}\label{fig:blowups:plane}
\end{figure}

\paragraph{Normal crossing and the recipe.} A change of variables $w=g(u)$ turns the loss into $K\circ g$ and the prior measure $\varphi(w)\dd w$ into the weight $\varphi(g(u))\abs{\det g'(u)}\dd u$. The volume and the zeta function of \cref{ch:volume} depend only on this pair, and they are monomial computations when both members are monomials in the same coordinates.

\begin{definition}[Normal crossing]\label{def:blowups:nc}
Let $\kappa\ge0$ be analytic and $\omega\ge0$ measurable near a point $P$ of a $d$-dimensional analytic manifold. The pair $(\kappa,\omega)$ is \term{normal crossing at} $P$ if there are analytic coordinates $u$ centred at $P$ in which, near $P$,
\[
\kappa=u_1^{2k_1}\cdots u_d^{2k_d}\,a(u),\qquad\omega=\abs{u_1}^{h_1}\cdots\abs{u_d}^{h_d}\,b(u),
\]
with nonnegative integers $k_j,h_j$ and positive infinitely differentiable functions $a,b$. The \term{candidates} at $P$ are the numbers $(h_j+1)/(2k_j)$ for the $j$ with $k_j>0$.
\end{definition}

The exponents of $\kappa$ are even because $\kappa\ge0$, and a positive factor such as $a$ or $b$, called a \term{unit}, changes no exponent. The recipe, which \cref{sec:blowups:resolution} justifies, has four steps.
\begin{enumerate}
\item Pull the loss and the prior back to the charts of a blow-up, multiplying the prior by $\abs{\det g'}$.
\item Check normal crossing at every point of the exceptional divisor in every chart, not only at the centres.
\item Where the check fails, blow up that point, and repeat.
\item When every point is normal crossing, $\lambda$ is the smallest candidate, and $m$ the largest number of coordinates attaining it at one point.
\end{enumerate}
\Cref{thm:toolkit:hironaka} guarantees that the process can be completed for every analytic loss; in general it needs blow-ups along smooth submanifolds,\sidenote{In $\R^d$ the blow-up along $\{w_1=\dots=w_k=0\}$ has the charts $w_1=u_1$, $w_i=u_1u_i$ for $2\le i\le k$, the other coordinates unchanged, with Jacobian $\abs{u_1}^{k-1}$.} but for two variables point blow-ups suffice. We drop the factor $\frac12$ of the quartic loss, which changes neither $\lambda$ nor $m$, since the volume function of $cK$ is $V(\eps/c)$.

\begin{computation}{The quartic loss $a^2+b^4$ by two blow-ups}\label{comp:blowups:quartic}
\step Near the origin $K=a^2+b^4$ vanishes only at $0$, and its Hessian there is $\diag(2,0)$: neither the Morse lemma nor normal crossing. The prior is $\frac14$ on $[-1,1]^2$.
\step Blow up the origin. In chart I, $K=u^2(1+u^2v^4)$ with Jacobian $\abs u$. The factor $1+u^2v^4$ is a unit, so the pair is normal crossing with $k=(1,0)$, $h=(1,0)$, and the candidate along $E_1=\{u=0\}$ is $\frac{1+1}{2}=1$.
\step In chart II, $K=v^2(u^2+v^2)$ with Jacobian $\abs v$. At the points $(u_0,0)$ with $u_0\ne0$ the factor $u^2+v^2$ is a unit and the candidate is again $1$. At the centre, the direction of the $b$-axis, along which $K$ grows only like $b^4$, the factor vanishes: not normal crossing.
\step Blow up the centre of chart II. In chart II.1, $(u,v)=(s,st)$, so $(a,b)=(s^2t,st)$, $K=s^4t^2(1+t^2)$, and the Jacobian is $\abs v\cdot\abs s=s^2\abs t$: candidates $\frac{2+1}{4}=\frac34$ for $s$ and $\frac{1+1}2=1$ for $t$. In chart II.2, $(u,v)=(st,t)$, so $(a,b)=(st^2,t)$, $K=t^4(1+s^2)$, and the Jacobian is $t^2$: candidate $\frac34$.
\step The factors $1+t^2$ and $1+s^2$ are units everywhere, so every point of charts I, II.1 and II.2, which cover the twice blown-up plane, is normal crossing. The candidates are $1$ on $E_1$ and $\frac34$ on the new curve $E_2$, which is $\{s=0\}$ in II.1 and $\{t=0\}$ in II.2; where $E_1$ and $E_2$ cross they differ. Hence $\lambda=\frac34$ and $m=1$, against $\frac d2=1$ and half the Hessian rank, $\frac12$.
\step Check by scaling: $a=\eps^{1/2}x$, $b=\eps^{1/4}y$ maps $\{a^2+b^4<\eps\}$ onto $\{x^2+y^4<1\}$, of area $B(\frac14,\frac32)=3.4961$, and for $\eps\le1$ the set lies in the square, so $V(\eps)=0.87402\,\eps^{3/4}$ exactly.
\end{computation}

\paragraph{One chart and the constant.} Chart II.2 alone, $(a,b)=(st^2,t)$, maps $\{t\ne0\}$ one-to-one onto $\{b\ne0\}$, the plane without a null set. It is not proper, since $s=a/b^2$ is unbounded near the $a$-axis, but the change of variables formula needs no properness. If $t^4(1+s^2)<\eps\le1$, then $\abs a\le\sqrt\eps$ and $\abs b\le\eps^{1/4}$, so the square imposes nothing, and
\begin{equation}\label{eq:blowups:quarticvolume}
\begin{aligned}
V(\eps)&=\frac14\int_{\R}\int_{\R}\ind\bigl\{t^4(1+s^2)<\eps\bigr\}\,t^2\dd t\dd s\\
&=\frac14\int_\R\frac23\Bigl(\frac{\eps}{1+s^2}\Bigr)^{3/4}\dd s=\frac{B(\frac12,\frac14)}{6}\,\eps^{3/4},
\end{aligned}
\end{equation}
which is the same $0.87402\,\eps^{3/4}$. The parameter $s$ labels the parabola $a=sb^2$ along which a point approaches the origin, and $s\in\R\cup\{\infty\}$ runs over the exceptional curve $E_2$: the constant is an integral over $E_2$, whose density $(1+s^2)^{-3/4}$ says how much prior mass approaches the truth along each parabola. This chart is one chart of a \term{weighted blow-up}\index{blow-up!weighted}, the one that \cref{ch:newton} reads off the Newton polygon. For the model's loss $\frac12(a^2+b^4)$ the volume is $0.87402\,(2\eps)^{3/4}=1.4699\,\eps^{3/4}$, and the Abelian lemma (\cref{lem:volume:abelian}) gives $\bar F(n)=\frac34\log n-\log(1.4699\,\Gamma(\frac74))+o(1)=\frac34\log n-0.301+o(1)$, as in \cref{ch:freeenergy}.

\paragraph{Three more by hand.} The regular $a^2+b^2$ becomes $u^2(1+v^2)$ in chart I, with Jacobian $\abs u$: one blow-up, candidate $1=\frac d2$. Squaring halves the candidate: $(a^2+b^2)^2$ becomes $u^4(1+v^2)^2$, so $\lambda=\frac12$, and indeed $V(\eps)=\frac\pi4\eps^{1/2}$. The product model at a generic truth, $(ab-c)^2$ with $0<c<1$, needs no blow-up: near a point of the hyperbola $ab=c$, where $a\ne0$, the functions $(u,v)=(ab-c,a)$ are analytic coordinates with Jacobian determinant $-a$, in which $K=u^2$; so $\lambda=\frac12$, $m=1$, a case of \cref{prop:blowups:morsebott} below.

\paragraph{What one had to notice.} Each direction of approach to the origin sees a simpler function than the loss, and a blow-up makes the directions into coordinates. Along most directions the quartic loss grows like $a^2$; along the $b$-axis it grows like $b^4$, and that direction, blown up again, carries the learning coefficient. The candidate $\frac34$ combines the order $4$ of the loss along $E_2$ with the order $2$ of the Jacobian, which records how much area the two blow-ups squeeze near $E_2$; the Hessian sees only the first blow-up.

\section{The theorem: what a resolution computes}\label{sec:blowups:resolution}

Resolutions are not unique, but they all give the same answer, because they compute the volume and the zeta function, which do not refer to coordinates. The \term{learning coefficient} and the \term{real log canonical threshold}\index{real log canonical threshold} of $K$ with respect to $\varphi$ are one number, with $-\lambda$ the largest pole of $\zeta(z)=\int K^z\varphi$ (\cref{ch:volume}): the first name is Watanabe's, for its statistical role, the second the geometers'. The prior $\varphi\ge0$ is bounded and need not be normalized. The engine is the monomial computation of \cref{ch:volume} with a weight that may vanish, since a partition of unity will cut the weight into pieces.

\begin{lemma}[Monomials with a continuous weight]\label{lem:blowups:monomial}
Let $k,h\in\N^d$ with $k\ne0$, and let $c\colon[0,1]^d\to[0,\infty)$ be continuous. Put $r_j=(h_j+1)/(2k_j)$ when $k_j>0$ and $r_j=\infty$ when $k_j=0$, $\lambda=\min_jr_j$, $J=\{j:r_j=\lambda\}$ and $m=\abs J$. Write $u=(u_J,u')$, with $u'$ the coordinates outside $J$, and $u^h=\prod_ju_j^{h_j}$. Then the integral over the face $\{u_J=0\}$
\begin{equation}\label{eq:blowups:constant}
C_c=\frac{1}{(m-1)!\,\lambda\prod_{j\in J}2k_j}\int_{[0,1]^{d-m}}c(0_J,u')\prod_{j\notin J}u_j^{\,h_j-2k_j\lambda}\dd u'
\end{equation}
converges, and as $\eps\to0$
\[
\int_{[0,1]^d}\ind\{u^{2k}<\eps\}\,u^h\,c(u)\dd u=\bigl(C_c+o(1)\bigr)\,\eps^\lambda\Bigl(\log\frac1\eps\Bigr)^{m-1}.
\]
If $c$ is positive at some point of the face, then $C_c>0$.
\end{lemma}

It is \cref{prop:volume:monomial} restated for one orthant of the cube, with $K$ exactly a monomial and a weight that may vanish, which is the form the proof of the theorem uses.

\begin{proof}
Extend $c$ to $[-1,1]^d$ by $c(u)=c(\abs{u_1},\dots,\abs{u_d})$, a continuous function. The integrand $\ind\{u^{2k}<\eps\}\abs{u^h}c(u)$ is even in every coordinate, so its integral over $[-1,1]^d$ is $2^d$ times the integral in the lemma. \Cref{prop:volume:monomial}, with $B=[-1,1]^d$, the factor $a\equiv1$ and the weight $b=c\ge0$, gives the integral over $[-1,1]^d$ as $(C+o(1))\eps^\lambda(\log\frac1\eps)^{m-1}$ with a finite constant $C$, which is $2^dC_c$: its face integral runs over $[-1,1]^{d-m}$, which gives the factor $2^{d-m}$, and $\prod_{j\in J}k_j=2^{-m}\prod_{j\in J}2k_j$. If $c(0_J,u'_0)>0$, then $c(0_J,\cdot)>0$ near $u'_0$ in the face, where the integrand of $C_c$ is positive almost everywhere, so $C_c>0$.
\end{proof}

The theorem allows the prior to be cut off by the boundary of the parameter set; in the charts this is recorded by \term{coordinate orthants}, sets $\{u:\sigma_ju_j\ge0\text{ for all }j\}$ with signs $\sigma_j\in\{1,-1\}$.

\begin{theorem}[The learning coefficient on a resolution]\label{thm:blowups:pullback}
Let $O\subset\R^d$ be open, $K\colon O\to[0,\infty)$ real analytic, and $\varphi\colon O\to[0,\infty)$ bounded and measurable, zero outside a compact set $W\subset O$; let $W_0=\{w\in W:K(w)=0\}$. Let $M$ be a $d$-dimensional real analytic manifold and $g\colon M\to O$ a proper real analytic map that, for some closed null set $N\subset O$ with null preimage, is an analytic isomorphism from $M\setminus g^{-1}(N)$ onto $O\setminus N$. Suppose that every point $P\in g^{-1}(W_0)$ is the centre of a \emph{coordinate cube}: analytic coordinates $u$ near $P$, with $u(P)=0$ and with $[-1,1]^d$ inside their range, in which, for $u\in[-1,1]^d$,
\begin{equation}\label{eq:blowups:nc}
K(g(u))=u^{2k}a(u),\qquad\varphi(g(u))\,\abs{\det g'(u)}=\abs{u^h}\,b(u)\,\ind_R(u),
\end{equation}
with $k,h\in\N^d$, $a$ and $b$ positive and infinitely differentiable, and $R$ a union of closed coordinate orthants, possibly empty. Let $\lambda_P$ be the smallest candidate $(h_j+1)/(2k_j)$ of the cube and $m_P$ the number of indices $j$ attaining it, with $\lambda_P=\infty$ if $R$ is empty, and suppose that $\lambda_P<\infty$ for some $P$. Then the smallest $\lambda_P$ is attained; call it $\lambda$, and let $m$ be the largest $m_P$ with $\lambda_P=\lambda$.
\begin{enumerate}
\item (Pulling back.) For every measurable $F\colon[0,\infty)\to[0,\infty]$, $\int_WF(K(w))\varphi(w)\dd w=\int_MF(K(g(u)))\varphi(g(u))\abs{\det g'(u)}\dd u$; in particular $V$ and $\zeta$ can be computed on $M$.
\item (Volume.) $V(\eps)\sim C\eps^\lambda(\log\frac1\eps)^{m-1}$ as $\eps\to0$, for a constant $C>0$.
\item (Zeta function.) $\zeta$ extends to a meromorphic function on $\C$ whose poles are negative rational numbers; the largest is $-\lambda$, of order $m$.
\item (Local values.) For $w\in W_0$, let $\lambda(w)$ be the smallest $\lambda_P$ over $P\in g^{-1}(w)$, and $m(w)$ the largest $m_P$ among those $P$ with $\lambda_P=\lambda(w)$. If $\lambda(w)<\infty$, both are attained, and for all small $r>0$, $V_{w,r}(\eps)=\int_{\{K<\eps\}\cap B(w,r)}\varphi\asymp\eps^{\lambda(w)}(\log\frac1\eps)^{m(w)-1}$.
\end{enumerate}
In particular, $\lambda$ and $m$ depend neither on the resolution nor on the cubes.
\end{theorem}

The integral over $M$ is computed chart by chart, $\abs{\det g'(u)}\dd u$ being the pull-back of Lebesgue measure. In \cref{comp:blowups:quartic}, $M$ is the twice blown-up plane, $N$ the origin, and the cubes are small squares around the points of $E_1\cup E_2$.

\begin{proof}
The plan is to change variables, cover $g^{-1}(W_0)$ by finitely many cubes in which $K\circ g$ is exactly a monomial, cut the integrals into pieces with a partition of unity, apply \cref{lem:blowups:monomial} to each piece, and recover the zeta function from the volume.

\emph{Step 1: pulling back.} On $M\setminus g^{-1}(N)$ the map $g$ is a diffeomorphism onto $O\setminus N$, so by the change of variables formula in charts, $\int_{O\setminus N}G(w)\dd w=\int_{M\setminus g^{-1}(N)}G(g(u))\abs{\det g'(u)}\dd u$ for every measurable $G\ge0$. Both $N$ and $g^{-1}(N)$ are null; take $G=F(K)\varphi$.

\emph{Step 2: exact monomials and finitely many cubes.} Let $P\in g^{-1}(W_0)$, with cube coordinates $u$ and, say, $k_1>0$. The functions $x_1=u_1a(u)^{1/(2k_1)}$ and $x_j=u_j$ for $j\ge2$ have a diagonal Jacobian matrix with positive entries at $u=0$, so they are coordinates near $P$, by the inverse function theorem. In them $K\circ g=x^{2k}$; $x_1$ has the sign of $u_1$, so $R$ is unchanged; and the weight is $\abs{x^h}\,\tilde b(x)\,\ind_R(x)$ with $\tilde b=b\,a^{-h_1/(2k_1)}\abs{\det(\partial u/\partial x)}$ positive and smooth. So on a small cube $Q_P=\{\abs{x_j}\le\delta_P\text{ for all }j\}$ the pair satisfies \eqref{eq:blowups:nc} with $a\equiv1$ and the same $k$, $h$, $R$. Let $Q_P'$ be the open cube of half the size. The set $g^{-1}(W_0)$ is compact, because $g$ is proper, so finitely many $Q'_{P_1},\dots,Q'_{P_N}$ cover it, and on the compact set $g^{-1}(W)\setminus\bigcup_iQ'_{P_i}$ the function $K\circ g$ is at least some $\delta_0>0$. Choose smooth $\psi_i\ge0$ on $M$, supported in $Q_{P_i}$, with $\sum_i\psi_i\le1$ everywhere and $\sum_i\psi_i=1$ on the compact set $\{u\in g^{-1}(W):K(g(u))\le\delta_0/2\}$: bump functions equal to $1$ on the half-cubes, divided by their sum and multiplied by a smooth cutoff that equals $1$ on that set and vanishes outside the union of the half-cubes, which contains the set.

\emph{Step 3: the volume.} Let $\omega=(\varphi\circ g)\abs{\det g'}$. For $\eps<\delta_0/2$, (i) gives $V(\eps)=\sum_iV_i(\eps)$ with $V_i(\eps)=\int_{Q_{P_i}}\ind\{x^{2k}<\eps\}\,\abs{x^h}\,\tilde b\,\psi_i\,\ind_R\dd x$ in the coordinates and exponents of the $i$-th cube. Split $Q_{P_i}$ into its closed orthants, keep those in $R$, and reflect and rescale each onto $[0,1]^d$, which multiplies $x^{2k}$ and the weight by constants. By \cref{lem:blowups:monomial}, $V_i(\eps)=(C_i+o(1))\eps^{\lambda_i}(\log\frac1\eps)^{m_i-1}$ with $C_i\ge0$ and $(\lambda_i,m_i)=(\lambda_{P_i},m_{P_i})$. Let $\lambda_*$ be the smallest $\lambda_i$ and $m_*$ the largest $m_i$ with $\lambda_i=\lambda_*$; then $V(\eps)=(C_*+o(1))\eps^{\lambda_*}(\log\frac1\eps)^{m_*-1}$, with $C_*\ge0$ the sum of the $C_i$ over the $i$ with $(\lambda_i,m_i)=(\lambda_*,m_*)$. Conversely, let $P\in g^{-1}(W_0)$ have $R\ne\emptyset$. The integrand of (i) is nonnegative, so $V(\eps)$ is at least its integral over the cube of $P$, which by \cref{lem:blowups:monomial}, applied to one orthant of $R$ with the weight $\tilde b>0$, is at least $c_P\eps^{\lambda_P}(\log\frac1\eps)^{m_P-1}$ with $c_P>0$. Comparing the bounds as $\eps\to0$, every such $P$ has $\lambda_P>\lambda_*$, or $\lambda_P=\lambda_*$ and $m_P\le m_*$. Since $(\lambda_*,m_*)$ is attained at one of the $P_i$, the minimum $\lambda$ is attained, $(\lambda,m)=(\lambda_*,m_*)$, and $C_*\ge c_{P_i}>0$. This proves (ii).

\emph{Step 4: the zeta function.} For $\Re z>0$, $\zeta(z)=\sum_i\int(x^{2k})^z\abs{x^h}\tilde b\,\psi_i\ind_R\dd x+E(z)$, where $E(z)=\int(1-\sum_i\psi_i)(K\circ g)^z\omega$. On the support of that integrand $K\circ g$ lies between $\delta_0/2$ and $\max K$, so $E$ is entire. After the same splitting into orthants, each other piece is a finite sum of integrals $\int_{[0,1]^d}\prod_jy_j^{2k_jz+h_j}\,c(y)\dd y$ with $c$ infinitely differentiable on $[0,1]^d$, which are meromorphic on $\C$ with poles among the numbers $-(h_j+1+\ell)/(2k_j)$, $\ell\in\N$, by \cref{def:toolkit:mellin}(ii) applied in one coordinate at a time. So $\zeta$ is meromorphic on $\C$, its poles are negative rational numbers, and by Step 3 none lies to the right of $-\lambda$. To see that $-\lambda$ is a pole of order exactly $m$, write $K^z=\abs z\int_K^\infty t^{z-1}\dd t$ for real $z<0$; by Tonelli's theorem, for $-\lambda<z<0$,
\[
\zeta(z)=\abs z\int_0^\infty t^{z-1}V(t)\dd t .
\]
The part over $t\ge\delta$ is at most $V(\infty)\delta^z$, bounded near $z=-\lambda$. On $t<\delta$, (ii) gives $V(t)=(C+\eta(t))t^\lambda(\log\frac1t)^{m-1}$ with $\eta(t)\to0$, and the substitution $t=e^{-s}$ gives $\int_0^\delta t^{\sigma-1}(\log\frac1t)^{m-1}\dd t=\Gamma(m)\sigma^{-m}+O(1)$ as $\sigma=z+\lambda\downarrow0$. Letting $\delta\to0$, $(z+\lambda)^m\zeta(z)\to\lambda\,\Gamma(m)\,C>0$ as $z\downarrow-\lambda$, so the pole at $-\lambda$ has order exactly $m$.

\emph{Step 5: local values.} Let $w\in W_0$ with $\lambda(w)<\infty$. The fibre $g^{-1}(w)$ is compact; cover it by finitely many half-cubes $Q'_P$, $P\in g^{-1}(w)$, with union $U$. Since $g$ is proper, $g^{-1}(\bar B(w,r))\subset U$ for small $r$; otherwise points outside $U$ with images tending to $w$ would accumulate at a point of $g^{-1}(w)\setminus U$, which is empty. The argument of Step 3, with functions $\psi_i$ for these finitely many cubes whose sum is $1$ on $g^{-1}(\bar B(w,r))$, bounds $V_{w,r}(\eps)$ above by a multiple of $\eps^{\lambda'}(\log\frac1\eps)^{m'-1}$ for the best pair $(\lambda',m')$ of the cubes; and for every $P\in g^{-1}(w)$ a small cube around $P$ maps into $B(w,r)$, so $V_{w,r}(\eps)\ge c\,\eps^{\lambda_P}(\log\frac1\eps)^{m_P-1}$. As in Step 3, $(\lambda',m')=(\lambda(w),m(w))$, both attained. Finally, (ii) and (iii) determine $\lambda$ and $m$ through $V$ and $\zeta$, which do not refer to $g$.
\end{proof}

\paragraph{Where resolutions come from.} Hironaka's theorem supplies them.\sidenote{Hironaka's two papers fill $217$ pages of the \emph{Annals of Mathematics} \parencite{Hironaka1964a,Hironaka1964b}.} In the form of \cref{thm:toolkit:hironaka}, every analytic $f$ with $f(w^*)=0$ has, near $w^*$, a proper analytic map $g$ that is an isomorphism off the zero set and makes $f\circ g$ and $\det g'$ normal crossing in the same coordinates. Applied to $K$, or to $K\varphi_1$ when the prior is an analytic factor $\varphi_1\ge0$ times a smooth positive one, it gives the hypotheses of \cref{thm:blowups:pullback} near each interior true parameter; resolving the defining functions of $W$ together with $K$ makes the region a union of orthants in each chart. \Textcite[Theorem~1]{Watanabe2013} states the result over a neighbourhood of the true parameters as his \term{standard representation}\index{standard representation}, $K(g(u))=u^{2k}$ and $\varphi(g(u))\abs{g'(u)}=b(u)\abs{u^h}$ with $b>0$, which is \eqref{eq:blowups:nc} with $a\equiv1$; he states that its proof is given in Theorem~6.1 of his 2009 book. With it, every analytic $K$ with such a prior has a learning coefficient, $\lambda$ is a positive rational number, and $1\le m\le d$. The black box is the existence of $g$.

\paragraph{When one blow-up is not enough.} A loss that vanishes on a curve with a singular point needs more blow-ups. The \term{cusp}\index{cusp} $a^2=b^3$ is the simplest singular curve, and its loss $(a^2-b^3)^2$ is the square of a function that changes sign, like the loss of a realizable regression.

\begin{computation}{The cusp $(a^2-b^3)^2$ by three blow-ups}\label{comp:blowups:cusp}
\step Write $K=f^2$ with $f=a^2-b^3$, and blow up the origin. In chart I, $f=u^2(1-uv^3)$ and $K=u^4\times$\,(a unit) near $E_1=\{u=0\}$, with Jacobian $\abs u$: candidate $\frac{1+1}4=\frac12$. In chart II, $f=v^2(u^2-v)$, with Jacobian $\abs v$; away from the centre the candidate along $E_1$ is again $\frac12$. At the centre, the curve $v=u^2$, the cusp seen in chart II and called its \term{strict transform}\index{strict transform}, is tangent to $E_1$: not normal crossing.
\step Blow up the centre of chart II. In chart II.2, $(u,v)=(st,t)$, $f=t^3(s^2t-1)$ and the Jacobian is $t^2$, so $K=t^6\times$\,(a unit) with candidate $\frac{2+1}6=\frac12$ along $E_2=\{t=0\}$. In chart II.1, $(u,v)=(s,st)$, $f=s^3t^2(s-t)$, the Jacobian is $s^2\abs t$, and $K=s^6t^4(s-t)^2$: three curves pass through the centre, $E_2=\{s=0\}$, $E_1=\{t=0\}$ and the strict transform $\{s=t\}$, and three curves through a point cannot all be coordinate axes.
\step Blow up the centre of chart II.1. In the chart $(s,t)=(p,pq)$, $K=p^{12}q^4(q-1)^2$ with Jacobian $p^4\abs q$; in the chart $(s,t)=(pq,q)$, $K=p^6q^{12}(p-1)^2$ with Jacobian $p^2q^4$. The new curve $E_3$, $\{p=0\}$ and $\{q=0\}$ respectively, has candidate $\frac{4+1}{12}=\frac5{12}$. It meets $E_1$ at $q=0$, the strict transform at $q=1$ and $E_2$ at $p=0$ in the second chart, each time transversally; the other curves have candidates $\frac12$, $\frac{0+1}2$ and $\frac12$.
\step Every point is now normal crossing, and $\lambda=\frac5{12}$, $m=1$: the singular point lowers $\lambda$ below the value $\frac12$ of the smooth points of the curve.
\step The constant. In the first chart, $p^{12}q^4(q-1)^2<\eps$ for $\abs p<(\eps/(q^4(q-1)^2))^{1/12}$, and over these $p$ the weight $\frac14p^4\abs q$ integrates to $\frac1{10}\abs q\,(\eps/(q^4(q-1)^2))^{5/12}$. The second chart covers the rest of $E_3$ through $q=1/p$, and
\[
\begin{aligned}
V(\eps)&\sim\frac{\eps^{5/12}}{10}\int_\R\abs q^{-2/3}\abs{1-q}^{-5/6}\dd q\\
&=\frac{B(\frac13,\frac16)+B(\frac12,\frac16)+B(\frac13,\frac12)}{10}\,\eps^{5/12}=1.9906\,\eps^{5/12}.
\end{aligned}
\]
\end{computation}

As for the quartic, the constant is an integral over the exceptional curve that carries $\lambda$, with the exponents $1-4\cdot\frac5{12}=-\frac23$ and $0-2\cdot\frac5{12}=-\frac56$ of \eqref{eq:blowups:constant} where other curves cross it. Scaling confirms it: $a=s^{1/2}x$, $b=s^{1/3}y$ with $s=\sqrt\eps$ turns $\{\abs{a^2-b^3}<s\}$ into $\{\abs{x^2-y^3}<1\}$, of area $7.9622=4\times1.9906$.

\paragraph{Real, not complex.} Algebraic geometry usually studies the \term{log canonical threshold}\index{log canonical threshold} over the complex numbers, and it can differ from the learning coefficient. The loss $(a^2+b^2+c^2)^2$ on $[-1,1]^3$ becomes $u^4(1+v^2+w^2)^2$ in the chart $(a,b,c)=(u,uv,uw)$, with Jacobian $u^2$, so $\lambda=\frac34$, and indeed $V(\eps)=\frac18\cdot\frac{4\pi}3\eps^{3/4}$. Over the complex numbers $1+v^2+w^2$ vanishes on a smooth surface, which gives the candidate $\frac{0+1}2=\frac12$; it has no real points, and the volume never sees it. The real threshold is at least the complex one \parencite[Corollary~2]{Saito2007}, so a complex computation only bounds $\lambda$ from below.

\section{The theorem: rules of the trade}\label{sec:blowups:rules}

Most learning coefficients in the literature come not from full resolutions but from rules that compare, add and multiply losses with known coefficients. The rules are statements about volumes and hold beyond analytic losses. Say that $(K,\varphi)$ has \term{learning coefficient $\lambda$ and multiplicity $m$ in the volume sense} if there are constants $0<c\le C$ with
\[
c\,\eps^\lambda\Bigl(\log\frac1\eps\Bigr)^{m-1}\le V(\eps)\le C\,\eps^\lambda\Bigl(\log\frac1\eps\Bigr)^{m-1}\qquad\text{for all small }\eps>0 .
\]
The pair $(\lambda,m)$ is then unique. For an analytic $K$ with a resolution, \cref{thm:blowups:pullback} gives these bounds with the pole data of $\zeta$; for non-analytic losses the volume sense is the one that makes sense: the tanh neuron with Gaussian inputs, whose $K$ is not analytic, has $K\asymp a^2b^2$ on its prior box and so, by rule (ii) below, the product model's pair. Write $(\lambda,m)\preceq(\lambda',m')$ if $\lambda<\lambda'$, or $\lambda=\lambda'$ and $m\ge m'$: the smaller pair has the larger volume.

\begin{proposition}[Rules of the trade]\label{prop:blowups:rules}
Let $\varphi,\varphi_1,\varphi_2$ be bounded priors vanishing outside a compact set $W$, and let the losses below be nonnegative on $W$.
\begin{enumerate}
\item (Invariance.) If $\Phi$ is an analytic isomorphism from an open set $O'$ onto an open set $O\supset W$, then $(K\circ\Phi,(\varphi\circ\Phi)\abs{\det\Phi'})$ and $(K,\varphi)$ have the same volume and zeta functions, hence the same $(\lambda,m)$.
\item (Comparison.) If $(K_1,\varphi)$ and $(K_2,\varphi)$ have $(\lambda_1,m_1)$ and $(\lambda_2,m_2)$ and $K_1\le c\,K_2$, then $(\lambda_1,m_1)\preceq(\lambda_2,m_2)$. If $(K_1,\varphi)$ has $(\lambda,m)$ and $c_1K_1\le K_2\le c_2K_1$ with $c_1,c_2>0$, then $(K_2,\varphi)$ has $(\lambda,m)$. If $(K,\varphi_1)$ has $(\lambda,m)$ and $c_1\varphi_1\le\varphi_2\le c_2\varphi_1$, then $(K,\varphi_2)$ has $(\lambda,m)$.
\item (Sums.) Let $w=(w',w'')\in\R^{d'}\times\R^{d''}$, $K(w)=K'(w')+K''(w'')$ and $\varphi(w)=\varphi'(w')\varphi''(w'')$. If $(K',\varphi')$ and $(K'',\varphi'')$ have $(\lambda',m')$ and $(\lambda'',m'')$, then $(K,\varphi)$ has $(\lambda'+\lambda'',\,m'+m''-1)$.
\item (Products.) With the same notation and $K=K'K''$, suppose the zeta functions of $(K',\varphi')$ and $(K'',\varphi'')$ are meromorphic on $\C$ with real poles, largest poles $-\lambda'$ and $-\lambda''$, of orders $m'$ and $m''$, and that their defining integrals converge for real $z>-\lambda'$ and $z>-\lambda''$ respectively. Then the zeta function of $(K,\varphi)$ is meromorphic on $\C$ with real poles, with largest pole $-\min(\lambda',\lambda'')$ of order $m'$ if $\lambda'<\lambda''$, $m''$ if $\lambda''<\lambda'$, and $m'+m''$ if $\lambda'=\lambda''$.
\item ($\lambda\le\frac d2$.) If $w_0$ is an interior point of $W$ with $K(w_0)=0$, $K$ is twice continuously differentiable near $w_0$, and $\varphi\ge c>0$ near $w_0$, then $V(\eps)\ge c'\eps^{d/2}$ for small $\eps$; so $\lambda\le\frac d2$ whenever $(K,\varphi)$ has a learning coefficient.
\item (Global is the minimum of local.) Let $K$ be continuous on $W$, $W_0=\{w\in W:K(w)=0\}$ nonempty, and suppose every $w\in W_0$ has a radius $r_w>0$ such that $(K,\varphi\ind_{B(w,r_w)})$ has a pair $(\lambda(w),m(w))$. Then the smallest $\lambda(w)$ is attained; call it $\lambda$; and $(K,\varphi)$ has $(\lambda,m)$ with $m$ the largest $m(w)$ for which $\lambda(w)=\lambda$.
\end{enumerate}
\end{proposition}

By \cref{thm:blowups:pullback}(iv), the hypothesis of (vi) holds for every analytic $K$ with a resolution, with $\lambda(w)$ and $m(w)$ the \term{local learning coefficient}\index{learning coefficient!local} and multiplicity at $w$, which \cref{ch:llc} estimates at trained networks.

\begin{proof}
The plan is to bound volume functions above and below by one another; only (iv) uses the zeta function. Throughout, $L=\log\frac1\eps$.

(i) By the change of variables formula, both pairs give the same $\int F(K)\varphi$ for every measurable $F\ge0$, in particular for $F=\ind_{[0,\eps)}$ and $F(t)=t^z$.

(ii) If $K_1\le cK_2$, then $\{K_2<\eps\}\subset\{K_1<c\eps\}$, so for small $\eps$
\[
c'_2\eps^{\lambda_2}L^{m_2-1}\le V_2(\eps)\le V_1(c\eps)\le C''\eps^{\lambda_1}L^{m_1-1},
\]
with constants from the two definitions; as $\eps\to0$ this forces $\lambda_1<\lambda_2$, or $\lambda_1=\lambda_2$ and $m_1\ge m_2$. If $c_1K_1\le K_2\le c_2K_1$, the same inclusions give $V_1(\eps/c_2)\le V_2(\eps)\le V_1(\eps/c_1)$, both of the order $\eps^\lambda L^{m-1}$. For the priors, $c_1V_{\varphi_1}\le V_{\varphi_2}\le c_2V_{\varphi_1}$.

(iii) Since $K',K''\ge0$, the set $\{K<\eps\}$ contains $\{K'<\eps/2\}\times\{K''<\eps/2\}$ and is contained in $\{K'<\eps\}\times\{K''<\eps\}$. By Fubini's theorem $V'(\frac\eps2)V''(\frac\eps2)\le V(\eps)\le V'(\eps)V''(\eps)$, and both bounds are of the order $\eps^{\lambda'+\lambda''}L^{(m'-1)+(m''-1)}$.

(iv) By Tonelli's theorem $\zeta(z)=\zeta'(z)\zeta''(z)$ for $\Re z>0$, and the product continues meromorphically with real poles. For real $z>-\lambda''$ the integral defining $\zeta''$ converges, so $\zeta''(z)=\int K''^z\varphi''>0$; for analytic losses with a resolution this convergence is Step~4 of the proof of \cref{thm:blowups:pullback}. If $\lambda'<\lambda''$, then $\zeta''$ is holomorphic and positive at $-\lambda'$, so the product has a pole of order $m'$ there and none to its right. If $\lambda'=\lambda''$, the leading Laurent terms multiply and the orders add.

(v) Since $K\ge0$ is minimal at the interior point $w_0$, $\nabla K(w_0)=0$, and Taylor's formula gives $K(w)\le M\abs{w-w_0}^2$ on a ball $B(w_0,\rho)$ where $\varphi\ge c$, with $M$ half the largest norm of the Hessian there. So $\{K<\eps\}$ contains the ball of radius $\sqrt{\eps/M}$ about $w_0$ once $\eps<M\rho^2$, and $V(\eps)\ge c\,v_d(\eps/M)^{d/2}$, with $v_d$ the volume of the unit ball. If also $V(\eps)\le C\eps^\lambda L^{m-1}$, then $\lambda\le\frac d2$.

(vi) The set $W_0$ is compact. Cover it by finitely many balls $B(w_i,r_{w_i})$; outside their union $K\ge\delta>0$ on $W$. Let $(\lambda_*,m_*)$ be the smallest of the pairs $(\lambda(w_i),m(w_i))$ for $\preceq$, and $V_w$ the volume of $(K,\varphi\ind_{B(w,r_w)})$. For $\eps<\delta$, $V(\eps)\le\sum_iV_{w_i}(\eps)\le C\eps^{\lambda_*}L^{m_*-1}$, and $V(\eps)\ge V_{w_i}(\eps)\ge c\,\eps^{\lambda_*}L^{m_*-1}$ for the $i$ attaining $(\lambda_*,m_*)$; so $(K,\varphi)$ has $(\lambda_*,m_*)$. For every $w\in W_0$, $V(\eps)\ge V_w(\eps)$, which as in (ii) gives $(\lambda_*,m_*)\preceq(\lambda(w),m(w))$; hence $\lambda_*$ is the smallest $\lambda(w)$ and $m_*$ the largest $m(w)$ among the points attaining it.
\end{proof}

\paragraph{Watanabe's example.} \Textcite{Watanabe2013} resolves $K=(ab+c)^2+a^2b^4$, ``the Kullback--Leibler distance of a neural network model'' from his 2009 book, with four charts, such as $(a,b,c)=(a_3,b_3,a_3b_3(b_3c_3-1))$, where $K=a_3^2b_3^4(c_3^2+1)$ with Jacobian $\abs{a_3b_3^2}$ and candidates $1$ and $\frac34$; he finds $\lambda=\frac34$, $m=1$. The rules give this in two lines. The shear $c'=ab+c$ has Jacobian $1$ and makes $K=c'^2+a^2b^4$, a sum in separate variables; on $\{K<\eps\}$ with $\eps\le0.1$, $\abs{ab}=\abs a^{1/2}(\abs ab^2)^{1/2}\le\eps^{1/4}$, so $\abs c<1$, and the cube in $(a,b,c')$ has the same small volumes as the cube in $(a,b,c)$. The loss $c'^2$ has $(\frac12,1)$, the monomial $a^2b^4$ has $(\frac14,1)$, and by (iii) $\lambda=\frac34$, $m=1$. For $\abs a\ge\sqrt\eps$ the substitution $b=\abs a^{-1/2}y$ gives $\{c'^2+a^2b^4<\eps\}$ the area $\abs a^{-1/2}B(\frac14,\frac32)\eps^{3/4}$, so with the prior $\frac18$, $V(\eps)\sim\frac12B(\frac14,\frac32)\eps^{3/4}=1.748\,\eps^{3/4}$. The local value is $\frac34$ all along the $a$-axis and $1$ along the $b$-axis away from the origin.

The last rule concerns the commonest shape of a set of true parameters: a smooth manifold along which the loss grows quadratically in every transverse direction. Following \textcite{WeiEtAl2023}, we call such a set \term{minimally singular}\index{minimally singular}; a function with such a minimum set is called Morse--Bott,\index{Morse--Bott function} after \textcite{Bott1954}. \Cref{ch:stages} applies the rule to the toy model of superposition, as \textcite{ChenEtAl2023} do.

\begin{proposition}[Minimally singular sets]\label{prop:blowups:morsebott}
Let $K\ge0$ be twice continuously differentiable near $w^*\in\R^d$, with $K(w^*)=0$. Suppose that near $w^*$ the zero set $\{K=0\}$ is a $C^2$ submanifold $S$ of dimension $s$, $0\le s<d$, and that $\nabla^2K(w^*)$ has rank $d-s$; equivalently, it is positive definite on the orthogonal complement of the tangent space $T_{w^*}S$, on which it vanishes. Let $\varphi$ be bounded, and bounded below by a positive constant near $w^*$. Then for every small $r>0$,
\[
V_{w^*,r}(\eps)=\int_{\{K<\eps\}\cap B(w^*,r)}\varphi\asymp\eps^{(d-s)/2}:
\]
the local learning coefficient at $w^*$ is $\frac{d-s}2$ and its multiplicity is $1$.
\end{proposition}

The case $s=0$ is the regular answer. The Hessian vanishes on $T_{w^*}S$ because $\nabla K\equiv0$ on $S$: differentiating along a curve $\gamma$ in $S$ gives $\nabla^2K\,\gamma'=0$; so the hypothesis is that the rank is as large as it can be. \Cref{thm:llc:hessian} bounds $\lambda(w^*)$ between $\frac r2$ and $\frac r2+\frac{d-r}4$ for a Hessian of rank $r$: minimally singular sets attain the lower bound, and the quartic model ($r=1$, $d=2$) the upper one.

\begin{proof}
The plan is to straighten $S$ and compare $K$ with the squared distance to it. Because $S$ is a $C^2$ submanifold, a $C^2$ diffeomorphism $\Phi$ from a neighbourhood of $0\in\R^{d-s}\times\R^s$ onto a neighbourhood of $w^*$, with $\Phi(0)=w^*$, maps $\{x=0\}$ onto $S$. Then $\tilde K=K\circ\Phi$ is $C^2$ and nonnegative, with $\tilde K(0,y)=0$ and, since each $(0,y)$ is a minimum, $\nabla_x\tilde K(0,y)=0$. Taylor's formula in $x$ gives
\[
\tilde K(x,y)=x^\top H(x,y)\,x,\qquad H(x,y)=\int_0^1(1-\tau)\,\nabla_x^2\tilde K(\tau x,y)\dd\tau,
\]
with $H$ continuous and $H(0,0)=\frac12\nabla_x^2\tilde K(0,0)$. At the critical point $w^*$ the Hessian transforms as a quadratic form, $\nabla^2\tilde K(0)=\Phi'(0)^\top\nabla^2K(w^*)\Phi'(0)$, of rank $d-s$. Its $y$-block vanishes because $\tilde K(0,y)\equiv0$, and its mixed block because $\nabla_x\tilde K(0,y)\equiv0$; so the $x$-block, positive semidefinite because $\tilde K\ge0$, has rank $d-s$ and is positive definite. By continuity there are $\rho>0$ and $0<c_1\le c_2$ with $c_1\abs x^2\le\tilde K(x,y)\le c_2\abs x^2$ on the box $Q_\rho=\{\abs x_\infty\le\rho,\ \abs y_\infty\le\rho\}$, where $\tilde\varphi=(\varphi\circ\Phi)\abs{\det\Phi'}$ is bounded above and below. For $\rho'\le\rho$ and small $\eps$, the set $\{x:c\abs x^2<\eps\}$ is a ball inside the cube of side $2\rho'$, of volume $v_{d-s}(\eps/c)^{(d-s)/2}$, so the volume of $(\tilde K,\tilde\varphi)$ on $Q_{\rho'}$ is $\asymp\eps^{(d-s)/2}$. Finally, for small $r$ there is $\rho'<\rho$ with $\Phi(Q_{\rho'})\subset B(w^*,r)\subset\Phi(Q_\rho)$, and by the change of variables formula $V_{w^*,r}$ lies between the volumes on $Q_{\rho'}$ and $Q_\rho$.
\end{proof}

\paragraph{The rules on the cast.} On the product model's cross away from the origin the Hessian of $a^2b^2$ has rank $1$ and the local pair is $(\frac12,1)$; at the origin it is the monomial value $(\frac12,2)$, the global pair by (vi). For rank-one regression at a point with $a\ne0$, $b=0$, the zero set $\{b=0\}$ has dimension $N$ and the Hessian is $\norm a^2I_M$ in $b$, so the local value is $\frac M2$, against the global $\frac{\min(M,N)}2$ (\cref{ch:catalogue,ch:llc}).

\section{Measuring it}\label{sec:blowups:measuring}

The script \texttt{05-blowups-check.py} checks the theorem on the quartic $a^2+b^4$, the cusp $(a^2-b^3)^2$ and Watanabe's $(ab+c)^2+a^2b^4$, with the uniform prior on $[-1,1]^d$, $\lambda=\frac34,\frac5{12},\frac34$, $m=1$ and the constants $C=0.87402$, $1.99056$, $1.74804$ computed above. It computes the volume at $27$ values of $\eps$ from $10^{-1}$ to $10^{-14}$ in three ways: by quadrature of the exact lengths of the sections of $\{K<\eps\}$ in the original coordinates, which uses nothing about the resolution; as the sum over the leaf charts of \cref{comp:blowups:quartic,comp:blowups:cusp}, each a unit square by \eqref{eq:blowups:twocharts} applied at every blow-up, with the Jacobian as weight; and by counting which of $10^7$ prior draws (seeds $5005$ to $5007$) have $K<\eps$, with Clopper--Pearson $95\%$ intervals, wherever at least $20$ draws fall in the set. It runs in under a minute.

The chart sums agree with the original volumes to a relative $1.7\cdot10^{-7}$ (quartic) and $7.8\cdot10^{-7}$ (cusp) at every $\eps$, and to $10^{-12}$ at most values of $\eps$.\sidenote{Near the corners $(\pm1,1)$, where the cusp leaves the square, the sections in charts I and II.2 change their number of intervals; adaptive quadrature that is not told so loses the sixth digit.} The Monte Carlo intervals contain the exact volume at $12$ of $13$, $26$ of $27$ and $14$ of $14$ values of $\eps$ (correlated, through shared draws). In \cref{fig:blowups:volume}(a) the slopes are the learning coefficients, not $\frac d2$. In panel (b) the quartic's ratio to the prediction is $1$, because $V(\eps)=C\eps^{3/4}$ exactly for $\eps\le1$; Watanabe's is $0.858$ at $\eps=10^{-2}$ and $0.9999$ at $10^{-14}$; the cusp's is $0.656$, $0.841$, $0.926$ and $0.966$ at $\eps=10^{-2},10^{-6},10^{-10},10^{-14}$. The cusp converges slowly because its next pole is close: the curves with candidate $\frac12$ give a correction of order $\eps^{1/2}$, relative order $\eps^{1/12}$, and the script finds $V(\eps)=C\eps^{5/12}-1.000\,\eps^{1/2}$ to four digits from $10^{-6}$ to $10^{-14}$. The deficit is the thin band around the curve $a^2=b^3$ that the square cuts off where the curve leaves it: in the scaled coordinates below \cref{comp:blowups:cusp} the band beyond $y=s^{-1/3}$ has area $4s^{1/6}$ to leading order, which removes $\frac14\cdot4\eps^{1/2}$ from $V$. Least-squares slopes of $\log V$ over $[10^{-14},10^{-6}]$ are $0.7500$, $0.4095$ and $0.7494$; local slopes between $10^{-13}$ and $10^{-14}$ are $0.7500$, $0.4134$ and $0.7500$.

\begin{figure}[t]
\begin{fullwidth}
\centering
\begin{tikzpicture}
\begin{groupplot}[group style={group size=2 by 1,horizontal sep=17mm},splot,
  width=0.48\linewidth,height=0.33\linewidth,xmode=log,xlabel={$\eps$},xmin=4e-15,xmax=0.25,
  xtick={1e-14,1e-11,1e-8,1e-5,1e-2}]
\nextgroupplot[ymode=log,ylabel={$V(\eps)$},ymin=1e-12,ymax=3,title={(a) volumes},
  legend pos=south east,legend style={font=\scriptsize},ytick={1e-12,1e-9,1e-6,1e-3,1}]
\addplot[regularpred] table[x=eps,y=quartreg] {figures/data/05-blowups-volume.dat};
\addplot[theory] table[x=eps,y=quartpred] {figures/data/05-blowups-volume.dat};
\addplot[measured,ebars] table[x=eps,y=V,y error minus=lo,y error plus=hi] {figures/data/05-blowups-mc-quart.dat};
\addplot[measured2,ebars] table[x=eps,y=V,y error minus=lo,y error plus=hi] {figures/data/05-blowups-mc-cusp.dat};
\addplot[measured,mark=triangle*,mark size=1.9pt,ebars] table[x=eps,y=V,y error minus=lo,y error plus=hi] {figures/data/05-blowups-mc-wat.dat};
\addplot[theory,forget plot] table[x=eps,y=cusppred] {figures/data/05-blowups-volume.dat};
\addplot[theory,forget plot] table[x=eps,y=watpred] {figures/data/05-blowups-volume.dat};
\addplot[regularpred,forget plot] table[x=eps,y=cuspreg] {figures/data/05-blowups-volume.dat};
\addplot[regularpred,forget plot] table[x=eps,y=watreg] {figures/data/05-blowups-volume.dat};
\legend{regular: slope $\frac d2$,singular: $C\eps^\lambda$,$a^2+b^4$ (Monte Carlo),$(a^2-b^3)^2$,$(ab+c)^2+a^2b^4$}
\nextgroupplot[ylabel={$V(\eps)/(C\eps^\lambda)$},ymin=0.55,ymax=1.06,title={(b) divided by the prediction},
  legend pos=south east,legend style={font=\scriptsize}]
\addplot[theory,domain=4e-15:0.25,samples=2] {1};
\addplot[ink,line width=0.6pt] table[x=eps,y=quartcomp] {figures/data/05-blowups-volume.dat};
\addplot[only marks,mark=o,mark size=1.6pt,color=ink,each nth point=2] table[x=eps,y expr=\thisrow{quartchart}/\thisrow{quartpred}] {figures/data/05-blowups-volume.dat};
\legend{singular prediction,exact (original coordinates),exact (charts)}
\addplot[ink,line width=0.6pt,forget plot] table[x=eps,y=cuspcomp] {figures/data/05-blowups-volume.dat};
\addplot[only marks,mark=o,mark size=1.6pt,color=ink,each nth point=2,forget plot] table[x=eps,y expr=\thisrow{cuspchart}/\thisrow{cusppred}] {figures/data/05-blowups-volume.dat};
\addplot[ink,line width=0.6pt,densely dashed,forget plot] table[x=eps,y=watcomp] {figures/data/05-blowups-volume.dat};
\node[font=\scriptsize,anchor=south west] at (axis cs:1e-14,1.003) {$a^2+b^4$};
\node[font=\scriptsize,anchor=north east] at (axis cs:1e-4,0.935) {$(ab+c)^2+a^2b^4$};
\node[font=\scriptsize,anchor=north east] at (axis cs:3e-9,0.88) {$(a^2-b^3)^2$};
\end{groupplot}
\end{tikzpicture}
\end{fullwidth}
\caption{The volume $V(\eps)$ of three losses with the uniform prior on $[-1,1]^d$. (a) Monte Carlo estimates from $10^7$ prior draws (seeds $5005$ to $5007$) with Clopper--Pearson $95\%$ intervals, where at least $20$ draws fall in the set; bronze: the prediction $C\eps^\lambda$ of \cref{thm:blowups:pullback} with $\lambda=\frac34,\frac5{12},\frac34$ and the constants of \cref{comp:blowups:quartic,comp:blowups:cusp} and Watanabe's example; slate dashed: slope $\frac d2$ through the value at $\eps=10^{-1}$. (b) The exact volumes divided by $C\eps^\lambda$: lines from sections in the original coordinates (dashed for Watanabe's example), open circles from the sums over the leaf charts, which agree with the lines to $10^{-6}$. Script: \texttt{figures/scripts/05-blowups-check.py}.}\label{fig:blowups:volume}
\end{figure}

The script also estimates $\zeta(z)=\E[K(w)^z]$ under the prior by Monte Carlo in the original coordinates, at $z>-\lambda/4$, where $K^z$ has a finite fourth moment,\sidenote{For $z\le-\lambda/2$ the variance of $K^z$ is infinite, and the sample mean, dominated by rare huge terms, is unreliable at every practical sample size.} and on the charts by quadrature in a form that shows the pole: for the quartic, charts II.1 and II.2 contribute $\frac1{4z+3}$ times $\int_0^1t^{2z+1}(1+t^2)^z\dd t$ and $\int_0^1(1+s^2)^z\dd s$. \Cref{tab:blowups:zeta} shows the two agreeing within the intervals, and $(z+\lambda)\zeta(z)$ approaching the residue $\lambda C$ of Step~4 of the proof; the residues computed on the charts, $0.65551439$ and $0.82939962$, equal $\lambda C$ to eight digits.

\begin{table}[t]
\begin{fullwidth}
\centering
\small
\begin{tabular}{@{}rccc@{\hspace{7mm}}rccc@{}}
\toprule
\multicolumn{4}{c}{$a^2+b^4$, $\lambda=\frac34$} & \multicolumn{4}{c}{$(a^2-b^3)^2$, $\lambda=\frac5{12}$}\\
\cmidrule(r{4mm}){1-4}\cmidrule{5-8}
$z$ & Monte Carlo & charts & $(z+\lambda)\zeta$ & $z$ & Monte Carlo & charts & $(z+\lambda)\zeta$\\
\midrule
$-0.05$ & $1.06103\pm0.00005$ & $1.06107$ & $0.7428$ & $-0.025$ & $1.07211\pm0.00005$ & $1.07214$ & $0.4199$\\
$-0.10$ & $1.13168\pm0.00011$ & $1.13177$ & $0.7357$ & $-0.05$ & $1.15564\pm0.00011$ & $1.15571$ & $0.4238$\\
$-0.15$ & $1.21432\pm0.00020$ & $1.21450$ & $0.7287$ & $-0.10$ & $1.36889\pm0.00032$ & $1.36913$ & $0.4336$\\
$-0.50$ & & $2.73545$ & $0.6839$ & $-0.30$ & & $4.58629$ & $0.5351$\\
$-0.70$ & & $13.21920$ & $0.6610$ & $-0.40$ & & $44.72631$ & $0.7454$\\
$-\frac34$ & & residue & $0.6555$ & $-\frac5{12}$ & & residue & $0.8294$\\
\bottomrule
\end{tabular}
\end{fullwidth}
\caption{The zeta function of two losses with the uniform prior on $[-1,1]^2$. Monte Carlo: mean of $K^z$ over $10^7$ prior draws (seed $5050$) with a Student $95\%$ interval, for $z>-\lambda/4$ only. Charts: quadrature on the leaf charts. Last columns: $(z+\lambda)\zeta(z)$ from the charts, tending to the residue $\lambda C$ at $-\lambda$; slowly for the cusp, whose next pole is at $-\frac12$. Script: \texttt{figures/scripts/05-blowups-check.py}.}\label{tab:blowups:zeta}
\end{table}

A resolution also says what the posterior looks like near a singularity. The script \texttt{05-blowups-posterior.py} computes the Gibbs measure $\pi_n\propto e^{-nK}\varphi$ of the quartic model's population loss $K=\frac12(a^2+b^4)$, the posterior with the fluctuation of the data removed, whose normalizer is the partition function of \cref{ch:volume}. In the coordinates $s=a/b^2$ and $E=nK$ of chart II.2, the computation behind \eqref{eq:blowups:quarticvolume} gives $\pi_n$ the density $\propto E^{\lambda-1}e^{-E}(1+s^2)^{-3/4}$ while $E<n/2$: the energy $nK$ is $\mathrm{Gamma}(\frac34,1)$ at every $n$, independent of the point $s$ of the exceptional curve where the parameter approaches the truth; a regular model with two parameters would have $\mathrm{Gamma}(1,1)$. The script does not use this: it computes $\pi_n$ by midpoint quadrature on a $1500\times1500$ grid in the original coordinates, on a box holding all but $e^{-60}$ of the mass, for $n=10^2,10^4,10^6$. In \cref{fig:blowups:posterior} the three distribution functions of $nK$ coincide with that of $\mathrm{Gamma}(\frac34)$ to within $0.0008$, and the mean of $nK$ is $0.7500$ at all three $n$, against the regular $1$. Among $2000$ draws at $n=10^4$ the Spearman correlation of $\abs\theta$, $\theta=\arctan s$, with $nK$ is $0.006$, and the quartiles of $\theta$ are $-1.187$, $-0.004$, $1.125$, against $\mp1.144$ and $0$ for the density $\propto\abs{\cos\theta}^{-1/2}$. With data the posterior acquires a random tilt in these coordinates (\cref{thm:toolkit:empirical,ch:freeenergy}); that $\Ewb{nK}$ is about $\lambda/\beta$ is \cref{prop:temperature:mean}.

\begin{figure}[t]
\begin{fullwidth}
\centering
\begin{tikzpicture}
\begin{groupplot}[group style={group size=3 by 1,horizontal sep=14mm},splot,width=0.34\linewidth,height=0.29\linewidth]
\nextgroupplot[xlabel={$a$},ylabel={$b$},xmin=-0.045,xmax=0.045,ymin=-0.32,ymax=0.32,
  xtick={-0.04,0,0.04},ytick={-0.3,0,0.3},title={(a) original coordinates}]
\addplot[draws,mark size=0.5pt] table[x=a,y=b] {figures/data/05-blowups-post-draws.dat};
\nextgroupplot[xlabel={$\theta=\arctan(a/b^2)$},ylabel={$nK$},ymode=log,xmin=-1.62,xmax=1.62,ymin=1e-4,ymax=12,
  xtick={-1.5708,0,1.5708},xticklabels={$-\frac\pi2$,$0$,$\frac\pi2$},ytick={1e-4,1e-2,1},title={(b) resolved coordinates}]
\addplot[draws,mark size=0.5pt] table[x=theta,y=E] {figures/data/05-blowups-post-draws.dat};
\nextgroupplot[xlabel={$x$},ylabel={$\pi_n(nK\le x)$},xmode=log,xmin=1e-3,xmax=9,ymin=0,ymax=1.02,
  xtick={1e-3,1e-2,1e-1,1},title={(c) the law of $nK$},legend pos=north west,legend style={font=\scriptsize,fill=none}]
\addplot[regularpred] table[x=x,y=reg] {figures/data/05-blowups-post-cdf.dat};
\addplot[theory] table[x=x,y=pred] {figures/data/05-blowups-post-cdf.dat};
\addplot[measured,each nth point=4] table[x=x,y=n2] {figures/data/05-blowups-post-cdf.dat};
\addplot[measured2,each nth point=4] table[x=x,y=n4] {figures/data/05-blowups-post-cdf.dat};
\addplot[measured3,each nth point=4] table[x=x,y=n6] {figures/data/05-blowups-post-cdf.dat};
\legend{regular: Gamma$(1)$,singular: Gamma$(\frac34)$,$n=10^2$,$n=10^4$,$n=10^6$}
\end{groupplot}
\end{tikzpicture}
\end{fullwidth}
\caption{The Gibbs measure $\pi_n\propto e^{-nK}\varphi$ of the quartic model's population loss $K=\frac12(a^2+b^4)$, uniform prior on $[-1,1]^2$, by midpoint quadrature on a $1500\times1500$ grid in the original coordinates. (a) $2000$ draws at $n=10^4$ (seed $5004$). (b) The same draws in the coordinates of chart II.2: the point $\theta$ of the exceptional curve and the energy $nK$, which are independent. (c) The distribution function of $nK$ at three sample sizes, against $\mathrm{Gamma}(\frac34)$ (bronze) and the regular $\mathrm{Gamma}(1)$ (slate dashed); the measured values differ from $\mathrm{Gamma}(\frac34)$ by at most $0.0008$. Script: \texttt{figures/scripts/05-blowups-posterior.py}.}\label{fig:blowups:posterior}
\end{figure}

\section{The ledger}\label{sec:blowups:ledger}

A learning coefficient is computed by changing coordinates until the loss and the prior are monomials. Blow-ups do it and change no integral, so $\lambda$ and $m$ are read off the exponents on the charts of a resolution, every resolution gives the same answer, and the constants of the volumes are integrals over exceptional curves. Most coefficients in practice come from rules: comparison, sums, products, minima of local values, and $\frac{d-s}2$ on minimally singular sets.

\ledgercard{the learning coefficient $\lambda$ and its multiplicity $m$, computed on a resolution}{$\frac d2$ and $m=1$: the Morse lemma makes a nondegenerate minimum a sum of squares, which one blow-up resolves}{$\lambda=\min(h_j+1)/(2k_j)$ over the charts of a resolution with $K\circ g=u^{2k}$ and $\varphi\abs{\det g'}=b\abs{u^h}$, $m$ the most ties at one point (\cref{thm:blowups:pullback}); the rules of \cref{prop:blowups:rules,prop:blowups:morsebott}}{the quartic model $K=\frac12(a^2+b^4)$: two blow-ups, $\lambda=\frac34$, $m=1$, $V(\eps)=1.4699\,\eps^{3/4}$, against $\frac d2=1$ and half the Hessian rank, $\frac12$; the cusp $(a^2-b^3)^2$: three blow-ups, $\frac5{12}$}{chart sums equal original volumes to $10^{-6}$; local slopes $0.7500$, $0.4134$, $0.7500$ at $\eps=10^{-14}$; under the Gibbs measure $\propto e^{-nK}\varphi$ the energy $nK$ has mean $0.7500$ and law $\mathrm{Gamma}(\frac34)$ to $0.0008$}{analyticity (Hironaka's theorem); a resolution in which the prior and the boundary of $W$ are normal crossing too; a prior bounded above and below near the true parameters; knowing $K$}

\section{Where the theory stops}\label{sec:blowups:stops}

\paragraph{What was assumed.} \Cref{thm:blowups:pullback} assumes an analytic loss and a resolution in hand in which the prior, cut off by the boundary of the parameter set, is normal crossing too. The boundary matters: the product model at the truth $ab=1$ on $[-1,1]^2$ has true parameters only at the corners $\pm(1,1)$, where $(ab-1)^2=(s+t-st)^2$, with $s=1-a\ge0$ and $t=1-b\ge0$, vanishes at one point of a quadrant; the blow-up $s=x$, $t=xy$ gives $x^2(1+y-xy)^2$ with Jacobian $x$, so $\lambda=1$, not the $\frac12$ of a hyperbola crossing the square (\cref{ex:blowups:corner}). The rules assume priors bounded above and below near the true parameters; a prior that vanishes there raises $\lambda$ (\cref{ex:blowups:prior}), and Jeffreys' prior gives $\lambda\ge\frac d2$ \parencite{Watanabe2024}. And all of it assumes that the population loss $K$ is known; in practice it is not, and \cref{ch:temperature,ch:llc} estimate $\lambda$ from data.

\paragraph{Where it breaks.} Analyticity is not a technicality. The loss $K(a)=e^{-1/a^2}$ is infinitely differentiable and vanishes only at $0$,\sidenote{All its derivatives vanish at $0$, so its Taylor series there is $0$: the standard smooth function that is not analytic.} but $\{K<\eps\}=\{\abs a<(\log\frac1\eps)^{-1/2}\}$: with the uniform prior on $[-1,1]$, $V(\eps)=(\log\frac1\eps)^{-1/2}$, no power of $\eps$, a zeta function that diverges for every $z<0$, and a population free energy that grows like $\frac12\log\log n$. No resolution exists for it; \textcite{Watanabe2010a} notes the corresponding failure for a prior factor $e^{-1/\norm w^2}$. ReLU networks are analytic only piecewise; the ReLU neuron escapes because its $K=\frac14a^2b^2$ happens to be a polynomial. Second, Hironaka's theorem is an existence theorem. Constructive proofs exist \parencite{Villamayor1989,BierstoneMilman1997}, but the learning coefficients known exactly come from resolutions found by hand for particular families (\cref{ch:catalogue}), such as reduced-rank regression \parencite{AoyagiWatanabe2005} and deep linear networks \parencite{Aoyagi2024}. Third, complex computations answer a different question: for $(a^2+b^2+c^2)^2$ the complex threshold is $\frac12$ and the learning coefficient $\frac34$, the value $\frac nd$ that \textcite[Proposition~3.3]{Saito2007} proves for a zero of order $d$ in $n$ variables whose leading form has no other real zero. Finally, the asymptotics can be far away: the cusp's volume is $3.4\%$ below its leading term at $\eps=10^{-14}$, and a slope fitted over the range $[10^{-6},10^{-1}]$ that Monte Carlo reaches easily is $0.385$, against $\lambda=0.417$.

\paragraph{Open problems.} Computing the learning coefficients of the networks people train is open beyond the families of \cref{ch:catalogue}: exact values are known for one tanh layer with a single input, but not for one tanh layer with several inputs and outputs and more than four hidden units, nor for networks with two nonlinear layers. We know of no practical algorithm for the real log canonical threshold in many variables. The complex one is minus the largest root of the Bernstein--Sato polynomial \parencite{Kollar1997}, which computer algebra finds in small examples, but it only bounds $\lambda$ from below. And a theory of learning coefficients for losses that are smooth or piecewise analytic but not analytic, which would cover ReLU networks without luck, would need a replacement for Hironaka's theorem; \cref{ch:llc} measures such networks anyway.

\begin{exercises}
\item \consolidation (The blow-up as a manifold.) Show that the charts \eqref{eq:blowups:charts} agree on their overlap through $(u,v)\mapsto(1/v,uv)$, that $g$ is proper, and that $E$ is a circle. Derive \eqref{eq:blowups:twocharts}.\label{ex:blowups:manifold}

\item \consolidation (The constant of the quartic.) Show that the residue of $\zeta$ at $-\frac34$ for $a^2+b^4$ is $\frac18\bigl(\int_{-1}^1(1+s^2)^{-3/4}\dd s+\int_{-1}^1\abs t^{-1/2}(1+t^2)^{-3/4}\dd t\bigr)$, that $t=1/s$ turns the second integral into $\int_{\abs s>1}(1+s^2)^{-3/4}\dd s$, and hence that $C=\frac16B(\frac12,\frac14)=\frac14B(\frac14,\frac32)$.\label{ex:blowups:constant}

\item \consolidation (Units and powers.) Show that $a^2b^2+a^4b^4$ has $(\lambda,m)=(\frac12,2)$, that $a^2b^2c^2$ has $(\frac12,3)$ with no blow-up (check it with the product rule), and with one blow-up that $(a^2+b^2)^k$ has $\lambda=\frac1k$ and $(a^2+b^2+c^2)^k$ has $\lambda=\frac3{2k}$.\label{ex:blowups:units}

\item \consolidation (The chain, point by point.) Use \cref{prop:blowups:morsebott} and the monomial rule to find the local pair at every true parameter of the chain $(a_1a_2a_3)^2$ on $[-1,1]^3$, boundary points included, and deduce the global pair $(\frac12,3)$ from rule (vi).\label{ex:blowups:chain}

\item \byhand (A loss with two branches.) Show that $K=a^2b^2+a^4=a^2(a^2+b^2)$ becomes $u^2v^4(1+u^2)$ in chart II, so that $\lambda=\frac12$ and $m=2$, the product model's pair, although $K$ is not a monomial. Check that $V(\eps)\sim\frac14\eps^{1/2}\log\frac1\eps$ with the uniform prior.\label{ex:blowups:twobranches}

\item \byhand (Tangent parabolas.) Resolve $K=(a^2-b^4)^2$, which vanishes on the parabolas $a=\pm b^2$, by two point blow-ups, and show that $\lambda=\frac38$, $m=1$. \hint{in the chart $(a,b)=(st^2,t)$, $K=t^8(s-1)^2(s+1)^2$ with Jacobian $t^2$.}\label{ex:blowups:parabolas}

\item \byhand\hard (A higher cusp.) Show that $K=(a^2-b^5)^2$ needs four point blow-ups, that on the last exceptional curve $K\sim p^{20}$ with Jacobian $\abs p^6$, and that $\lambda=\frac7{20}$, $m=1$. \hint{the first blow-up gives $f=v^2(u^2-v^3)$ in chart II, a cusp again, with an extra $v^2$ in $f$ and $v$ in the Jacobian. Check by scaling $a=s^{1/2}x$, $b=s^{1/5}y$.}\label{ex:blowups:highercusp}

\item \byhand\hard (A network-shaped loss.) Let $K=(ab-cd)^2$ near $0\in\R^4$, the loss of a sum of two products. Blow up the origin; in the chart $a=u$, $b=uv_2$, $c=uv_3$, $d=uv_4$ show that $K=u^4(v_2-v_3v_4)^2$ with Jacobian $\abs u^3$, that $v_2-v_3v_4$ can serve as a coordinate, and that $\lambda=\frac12$, $m=1$: the singular point does not lower the value at the smooth points of the zero set. What does the Newton polyhedron of $K$ predict (\cref{ch:newton}), and why does it fail here? \hint{the candidates are $\frac{3+1}4$ and $\frac12$.}\label{ex:blowups:fourvars}

\item \variation (A prior that vanishes.)\label{ex:blowups:prior} Give $a^2+b^4$ the prior proportional to $\abs b^{\,h}$, $h\ge0$, on $[-1,1]^2$. Redo \cref{comp:blowups:quartic} with the prior in the weight and show that $\lambda=\frac{3+h}4$, which exceeds $\frac d2=1$ for $h>1$; with the prior proportional to $\abs a^h$, show that $\lambda=\frac{3+2h}4$. Why does this not contradict rule (v)?

\item \variation (A corner.)\label{ex:blowups:corner} Show that for $K=(ab-1)^2$ with the uniform prior on $[-1,1]^2$, $V(\eps)\sim\frac14\eps$, so $\lambda=1$, while for $(ab-c)^2$ with $0<c<1$, $V(\eps)\sim\log(\frac1c)\,\eps^{1/2}$. Explain the difference with the coordinates $s=1-a$, $t=1-b$ and one blow-up.

\item \measurement (Two poles.) Fit $V(\eps)=A\eps^\lambda+B\eps^{1/2}$, with $A$, $B$, $\lambda$ free, to the cusp's exact volumes on $[10^{-14},10^{-6}]$. How close does $\lambda$ come to $\frac5{12}$? Repeat with the Monte Carlo volumes and say what limits the accuracy.\label{ex:blowups:twopoles}

\item \measurement (The posterior with data.) Change the posterior script to sample the posterior $\propto e^{-nK_n}\varphi$ of the quartic model for one data set at $n=10^4$, and plot the draws in the coordinates $(\theta,nK)$. Is $nK$ still independent of $\theta$? Compare the mean of $nK$ over $100$ data sets with $\frac34$, and relate the difference to the Gibbs generalization error of \cref{ch:generalization}.\label{ex:blowups:data}

\item \research (Resolution by machine.) Implement recursive point blow-ups for polynomials $K(a,b)\ge0$ with $K(0,0)=0$, tracking strict transforms and Jacobians, and test it on the chapter's examples and on $(a^2-b^{2k+1})^2$, $k\le5$. How does the number of charts grow? Find a loss for which the Newton-polygon shortcut of \cref{ch:newton} fails and the blow-ups do not.\label{ex:blowups:machine}
\end{exercises}

\notesandrefs
Resolution of singularities in characteristic zero is the theorem of \textcite{Hironaka1964a,Hironaka1964b}. \Textcite{BernsteinGelfand1969} and \textcite{Atiyah1970} used it to continue $\int\abs f^z\varphi$ meromorphically; Atiyah's real analytic local form is the version that \textcite[Theorem~3.3]{Lin2011} attributes to Atiyah and says \textcite{Watanabe2009} uses as his Theorem~2.3. In learning theory, \textcite{Watanabe2001a} gives an algorithm for $\lambda$ by resolution, and Watanabe's standard representation \parencite[Theorem~1]{Watanabe2013} is the form in which the later theory uses Hironaka's theorem; his review \parencite{Watanabe2024} restates it. Constructive resolution goes back to \textcite{Villamayor1989} and \textcite{BierstoneMilman1997}.

In algebraic statistics, \textcite[Chapter~5]{DrtonEtAl2009} and \textcite[Chapter~17]{Sullivant2018} treat Bayesian integrals of algebraic models, and \textcite{Lin2011,Lin2017} develops the rules of \cref{prop:blowups:rules} for ideals, among them comparison, sums and products, and the Jacobian-rank bound that \cref{ch:llc} uses (Propositions~3.4 and~3.7 and Corollary~4.15 of the paper, in the numbering of its arXiv version); \textcite[Section~3]{Watanabe2024} lists the first three without proof. \Textcite{Saito2007} proves that the real threshold is at least the complex one and that it can be larger; \textcite{Kollar1997} surveys the complex theory.

\Cref{thm:blowups:pullback} is the book's own presentation of standard arguments: the monomial volumes of \cref{prop:volume:monomial}, a partition of unity on the resolved space, and the Mellin continuation of \cref{def:toolkit:mellin}. The cusp's three blow-ups are its classical minimal embedded resolution; the constants as integrals over exceptional curves and the sum-rule shortcut for Watanabe's example, whose loss he attributes to Example~1.6 of the 2009 book, are ours. \Cref{prop:blowups:morsebott} is proved here by comparison with the squared distance to the critical set; nondegenerate critical manifolds are \textcite{Bott1954}'s, and the proposition is the case of equality in the Hessian lower bound of \cref{ch:llc}.

\paragraph{Translating the literature.} \Textcite{Watanabe2013} writes the standard representation as $K(g(u))=u^{2k}$ and $\varphi(g(u))\abs{g'(u)}=b(u)\abs{u^h}$, our \eqref{eq:blowups:nc} with $a\equiv1$ and $\abs{\det g'}$ for his $\abs{g'}$, and calls $\lambda$ the real log canonical threshold and $m$ its multiplicity. \Textcite{Lin2011,Lin2017} writes $\mathrm{RLCT}_\Omega(f;\varphi)$ for the pair (threshold, multiplicity) of a function $f$ with amplitude $\varphi$ on a region $\Omega$: for $f=K$ our $(\lambda,m)$, and for an ideal $\langle f_1,\dots,f_r\rangle$ the pair $(2\lambda,m)$ when $K\asymp\sum_if_i^2$. \Textcite{Saito2007} writes $\mathrm{rlct}(f)$ for the supremum of the $c$ with $\abs f^{-c}$ locally integrable: $\lambda$ for $f=K$, $2\lambda$ for $K=f^2$. Algebraic geometers write $\mathrm{lct}(f)$ for the complex analogue; \textcite{Kollar1997} writes it $c_0(f)$ and computes it on a resolution as $\min_j(a_j+1)/b_j$, with $b_j$ the order of $f$ along the divisor $E_j$ and $a_j$ its discrepancy: for $f=K$, our $2k_j$ and $h_j$. \Textcite{Watanabe2024} writes $\theta$, $\Theta$, $\pi(\theta)$ for our $w$, $W$, $\varphi(w)$.

\printchapterbib
\end{refsection}
